In computational studies of humanistic tasks such as storytelling and poetry, we often encounter datasets whose features frustrate standard NLP pipelines. In this talk, I'll draw on three settings where my collaborators and I have run into this friction. First, I'll discuss user–chatbot conversations like WildChat, where creative requests are long, detailed, and frequently revised, making these logs a rich but also complicated view of how people write with chatbots. Second, I'll describe our development of a fine-grained narrative annotation scheme for very large-scale, heterogeneous web text in pretraining data, where operationalizing a humanities concept is complicated by data diversity. Third, I'll turn to poetry, where the linearization systems used to build training corpora routinely flatten or delete formatting that poets use deliberately. Across these cases, I'll argue that what counts as noise depends on who is reading and for what purpose, and that taking humanistic perspectives seriously can help us build data practices, annotation schemes, and evaluations that better reflect how people actually use language.